\chapter{示例章节}

本章演示模板的常用环境：数学定理、复杂度记号与代码块。确认排版无误后，可以删除本文件，并在 \texttt{main.tex} 中移除对应的 \verb|\include| 行。

\section{数学环境}

\begin{definition}[同余]
    若 $m \mid (a - b)$，则称 $a$ 与 $b$ 模 $m$ 同余，记作 $a \equiv b \pmod{m}$。
\end{definition}

\begin{theorem}[费马小定理]
    设 $p$ 为素数，$a$ 为整数且 $p \nmid a$，则
    \[
        a^{p-1} \equiv 1 \pmod{p}.
    \]
\end{theorem}

\begin{example}
    取 $p = 7$，$a = 3$，则 $3^{6} = 729 \equiv 1 \pmod{7}$。
\end{example}

常用记号：$\mathbb{N}, \mathbb{Z}, \mathbb{Q}, \mathbb{R}$，$\lvert -x \rvert = x$（$x \ge 0$），$\lfloor \tfrac{7}{2} \rfloor = 3$，$\lceil \tfrac{7}{2} \rceil = 4$，时间复杂度 $O(n \log n)$。

\section{代码环境}

\begin{lstlisting}
#include <bits/stdc++.h>
using namespace std;

// 快速幂：O(log n)
long long qpow(long long a, long long b, long long mod) {
    long long res = 1 % mod;
    a %= mod;
    while (b > 0) {
        if (b & 1) res = res * a % mod;
        a = a * a % mod;
        b >>= 1;
    }
    return res;
}

int main() {
    cout << qpow(3, 6, 7) << '\n';  // 输出 1
    return 0;
}
\end{lstlisting}

\section{表格}

\begin{table}[h]
    \centering
    \begin{tabular}{lll}
        \toprule
        数据结构 & 查询 & 修改 \\
        \midrule
        树状数组 & $O(\log n)$ & $O(\log n)$ \\
        线段树   & $O(\log n)$ & $O(\log n)$ \\
        分块     & $O(\sqrt{n})$ & $O(1)$ \\
        \bottomrule
    \end{tabular}
    \caption{常见区间数据结构复杂度对比}
\end{table}
